\section{Artifact provenance and build instructions}\label{app:provenance}
\subsection{Reproducibility boundary}
Architecture source, runtime source, data authority, and experimental evidence have separate identities. The editable modular source is \code{main.tex} plus \code{sections/}; \code{ISINGFOLD\_RL\_ARCHITECTURE\_EN.tex} is its self-contained distribution form. Runtime checkpoints bind the source registry implemented by \code{checkpoint.py}. Prepared-v4 binds the independent CandidateBank export, corpus design, provenance bytes, public policy data, and partitioned target commitments. Target-bearing artifacts additionally bind an out-of-band publisher attestation v2, target-access receipt, and root-authenticated selected-partition certificate authority. Target-free plans bind only the public ground root. Evaluation reports bind raw receipts rather than relying on aggregate numbers alone.

The scientific compatibility chain is prepared schema v4, publisher attestation v2, and ground-certificate protocol v2 with a target-free root plus selected-partition receipts; selector-label manifest v4 and selector bundle v4; quality record v7 with shard/merge envelope v5 and quality preflight v2; release gate v6; representation selection v4 and RL-value freeze v4; checkpoint v3 and training-run receipt v4. Initializer plan, snapshot, bank manifest, and access artifacts use v1. Learned complete-system receipts use v3 with evaluation report v4 and three-seed aggregate v2. External tuning run and selection receipts use v2; external complete-system receipts use v3 with report v4; and the paired aggregate uses v4. Resumable evaluation uses plan v2, shard receipt v3, and merge receipt v2. The final-strength diagnostic uses config v1, sampler identity v1, plan v2, row v1, execution manifest v2, shard v1, merge v1, and audit receipt v2. The main experiment registry is \code{configs/rl\_grid\_v1.json}; the complete-system contract is \code{configs/complete\_system\_lac\_v1.json}; the validation-only stock tuning registry is \code{configs/external\_minorminer\_tuning\_v1.json}; the final-strength diagnostic is \code{configs/final\_strength\_audit\_v1.json}; and the isolated capacity diagnostic is \code{configs/rl\_capacity\_control\_v1.json}. The operational sequence and exact command lines are maintained in \code{documents/TRAINING\_OPERATIONS.md}.

\subsection{Compatibility rules}
\begin{enumerate}
\item Preserve original V2 contexts and labels. Create IF-Q3-S0 contexts for the new selector and logical tolerance; do not overwrite historical endpoint identifiers.
\item Require complete claim and group-action tensors for overlap trajectories. A disjoint static record is not a substitute.
\item Keep the frozen terminal selector separate from PPO critics and exclude ground-truth strength selection from deployable code paths.
\item Charge proposal work, reserve the terminal path, set RL dropout to zero, and store exact action supports.
\item Keep hidden witness and optimum-certificate records outside policy input; split by immutable base parent.
\item Plan and seal per-seed initializer banks before target-bearing training; preserve every rejected draw and fail closed on an exhausted conditional episode.
\item Shard evaluation by whole immutable base lineages, derive seeds from scientific identities rather than shard order, and accept only a complete nonoverlapping merge.
\item Treat any schema, source digest, publisher identity, ground root or partition, selector, initializer bank, support, budget, or seed-schedule change as a new artifact lineage.
\item Run Apollo workloads directly with the pinned virtual environment. Run Goose workloads only through Slurm allocations with the pinned Apptainer image.
\end{enumerate}

\subsection{How to compile and edit the delivered source}
The delivered single-file \code{ISINGFOLD\_RL\_ARCHITECTURE\_EN.tex} includes the style, figures, all sections and references inline. It requires no external diagram exports or bibliography database. Compile with:
\begin{Verbatim}
pdflatex -interaction=nonstopmode -halt-on-error ISINGFOLD_RL_ARCHITECTURE_EN.tex
pdflatex -interaction=nonstopmode -halt-on-error ISINGFOLD_RL_ARCHITECTURE_EN.tex
\end{Verbatim}
Run a third pass if page references change. Alternatively use \code{latexmk -pdf}. The source bundle also includes modular \code{main.tex} and the \code{sections/} directory, a Makefile, a source-digest manifest and a README.

\section{Worked checks for the mathematical contracts}\label{app:worked}
\subsection{A small planted logical optimum}
Take a four-cycle with $J_{12}=J_{23}=J_{34}=-1$ and $J_{41}=+1$, with zero fields. Any spin assignment leaves an odd number of preferred edge relations unsatisfied, so at least one edge contributes $+1$. The minimum energy is $-4+2=-2$, attained by all-plus. A gauge $s^\star=(+1,-1,+1,-1)$ changes the signs to $(+1,+1,+1,-1)$ and preserves $E_0=-2$. The planted solution is not unique. A witness embedding of this cycle proves that it can be placed on a host; it still gives no ordering among several valid embeddings' solve probabilities.

\subsection{Equal qubit count can have different cut certificates}
Consider two candidate locations for a three-qubit logical chain, one inducing a path and one a triangle, with the same absolute external-load vector $(2,0,2)$. Such loads can arise from a logical coupling of magnitude four allocated over two endpoint contacts. Both candidates use three qubits and total load $A=4$. The path's single-edge endpoint cut has $d/c=2$; the triangle's endpoint cut has $d/c=1$, and its zero-load middle-vertex cut has ratio zero. Thus their exact cut sufficient thresholds are 2 and 1 respectively.

At $F=1.5$, the triangle has positive margin on every cut, whereas the path fails this conservative sufficient condition. The lack of a path certificate does not prove a broken ground state or lower decoded quality. If its two external contacts have aligned drive, the absolute-load bound can be loose. If the program coupler range is one and each problem contact has magnitude two, the shrink-only rule gives scale $a=1/2$ for both candidates. Scaling changes the sampling landscape even though it preserves ideal energy ordering. This example justifies including geometry and allocation features, not assigning a quality reward from the cut threshold alone.

\subsection{Why a noisy four-strength maximum is optimistic}
Suppose every strength has true success probability $0.5$ and each receives one independent diagnostic read. The maximum of the four observed proportions is one unless all four reads fail, hence its expectation is $1-(1/2)^4=0.9375$, although the population maximum is $0.5$. Selecting one strength without those outcomes and estimating it on a fresh read has expectation $0.5$. Increasing read counts reduces but does not generally eliminate the maximum-selection bias.

\subsection{A trajectory is an RL sample, not eight supervised examples}
For a two-step search, action $a_0$ may temporarily increase conflict and enable a final high-quality rewrite $a_1$. A one-step reward based only on conflict reduction can prefer another action and change the objective. Terminal utility with $\gamma=1$ credits the chosen sequence. Counterfactual labels under a frozen continuation $\mu$ quantify only what $\mu$ does after $a_0$; they are not the best possible continuation. PPO can improve that continuation through on-policy updates if the candidate support contains the necessary next actions. It cannot create an action absent from that support.
